\documentclass[UTF8,11pt]{ctexart}
\usepackage[paperwidth=240mm,paperheight=200mm,margin=13mm,top=8mm]{geometry}
\usepackage{amsmath,amssymb,mathtools}
\usepackage{xcolor}
\pagestyle{empty}
\setlength{\parindent}{0pt}
\setlength{\parskip}{3pt}
\xeCJKDeclareCharClass{CJK}{"2460 -> "24FF}
\begin{document}
{\Large\bfseries 2026年考研数学二试题 答案与解析}\par\vspace{4pt}
\par\small\color{gray}\emph{选择题：1.A　2.B　3.A　4.D　5.C　6.B　7.D　8.B　9.A　10.D}\par
\par\small\color{gray}\emph{填空题：11.(0,2)　12.$\dfrac12$　13.$4$　14.$-2\pi$　15.$3\ln2-1$　16.$2$}\par
\par\small\color{gray}\emph{解答题：见下文。}\par
{\Large\bfseries 一、选择题（1～10）}\par\vspace{4pt}
\par\noindent\textbf{1. A. $a=\dfrac13$，$b=-1$}\par
$x\to0$ 时 $\sqrt[3]{1+x^2}-1\sim\dfrac13x^2$，$\arcsin x\sim x$。由 $f(x)=ax^2+bx+\arcsin x\sim ax^2+(b+1)x$ 与 $\dfrac{\sqrt[3]{1+x^2}-1}{f(x)}\to1$，需 $b+1=0$ 且 $a=\dfrac13$，故 $b=-1,\ a=\dfrac13$。
\par\noindent\textbf{2. B. $\lambda=\dfrac25$，$\mu=\dfrac15$}\par
令线性算子 $L$。由 $L(2\lambda y_1+\mu y_2)=f$ 得 $2\lambda+\mu=1$；由 $L(\lambda y_1-2\mu y_2)=0$ 得 $\lambda-2\mu=0$，即 $\lambda=2\mu$。联立得 $\mu=\dfrac15,\ \lambda=\dfrac25$。
\par\noindent\textbf{3. A. $\dfrac{\partial z}{\partial x}-\dfrac{\partial z}{\partial y}=\dfrac1a$}\par
设 $F=x-az-\mathrm{e}^{y+az}=0$，$F_x=1,\ F_y=-\mathrm{e}^{y+az},\ F_z=-a(1+\mathrm{e}^{y+az})$。则
$\dfrac{\partial z}{\partial x}=\dfrac{1}{a(1+\mathrm{e}^{y+az})}$，$\dfrac{\partial z}{\partial y}=-\dfrac{\mathrm{e}^{y+az}}{a(1+\mathrm{e}^{y+az})}$。相减得 $\dfrac{\partial z}{\partial x}-\dfrac{\partial z}{\partial y}=\dfrac{1+\mathrm{e}^{y+az}}{a(1+\mathrm{e}^{y+az})}=\dfrac1a$。
\par\noindent\textbf{4. D. $2\displaystyle\int_0^1\dfrac{Gm}{(x^2+1)^{\frac32}}\,\mathrm{d}x$}\par
对称性使水平分量抵消，竖直分量 $F=2\displaystyle\int_0^1\dfrac{Gm}{(x^2+1)}\cdot\dfrac{1}{\sqrt{x^2+1}}\,\mathrm{d}x=2\int_0^1\dfrac{Gm}{(x^2+1)^{\frac32}}\,\mathrm{d}x$。
\par\noindent\textbf{5. C}\par
若 $f$ 的图形在 $[-1,1]$ 凹，则分式函数 $\dfrac{f(x)-f(1)}{x-1}$ 在 $[-1,1)$ 单调递增。（A、B 仅在 $f$ 连续时成立；D 为 C 的逆命题，凹性不能仅由该分式单调推出，不充分。）
\par\noindent\textbf{6. B. $g(0)=1$，$g'(0)=\dfrac{2}{3}\mathrm{e}$}\par
$f(1)=\int_1^1=0$，故 $g(0)=1$。$f'(x)=\dfrac{\mathrm{e}^{x^3}}{1+x^6}\cdot3x^2$，$f'(1)=\dfrac{3\mathrm{e}}{2}$，故 $g'(0)=\dfrac{1}{f'(1)}=\dfrac{2}{3\mathrm{e}}$。
\par\noindent\textbf{7. D}\par
$f(x,y)=f(y,x)$ 对称，$\iint_D f=2\iint_{x\geqslant y}f$。在 $2n$ 等分下求和：$\iint_D f=\dfrac12\lim\limits_{n\to\infty}\sum\limits_{i=1}^{2n}\sum\limits_{j=1}^{i}f\left(\dfrac{i}{2n},\dfrac{j}{2n}\right)\dfrac1{n^2}$。
\par\noindent\textbf{8. B}\par
置换矩阵 $A$ 满足 $A^{-1}=A^{\mathrm{T}}$，而转置置换矩阵仍是置换矩阵，故 $A^{-1}$ 为置换矩阵。$A^*=|A|A^{-1}=\pm A^{-1}$，故 C、D 仅在 $|A|=\pm1$ 时成立，不恒真。
\par\noindent\textbf{9. A. $a=-1,\ b=-1$}\par
$A=\begin{pmatrix}1&0&1\\0&0&1\\1&1&3\end{pmatrix}$，$C=\begin{pmatrix}2&0\\1&1\\a&b\end{pmatrix}$。由 $AB=C$ 须 $C$ 的各列在 $A$ 的列空间中。对第一列 $Ax=(2,1,a)^{\mathrm{T}}$ 与第二列 $Ax=(0,1,b)^{\mathrm{T}}$ 分别消元，得 $a=-1,\ b=-1$ 时成立。
\par\noindent\textbf{10. D}\par
由 $AB+BA=A^2+B^2$ 得 $(A-B)^2=O$，故 $(A-B)^3=O$（A 对），$A-B$ 幂零只有零特征值（B 对）。若 $A,B$ 皆对角则 $AB=BA$，代入得 $A=B$，矛盾，故 C 对。而 $A-B$ 幂零指数为 $2$，其零特征值的几何重数 $\geqslant2$，故「只有一个线性无关特征向量」错误。答案 D。
{\Large\bfseries 二、填空题（11～16）}\par\vspace{4pt}
\par\noindent\textbf{11. $(0,2)$}\par
$x\to0^+$：$\arctan x\sim x$，积分近 $x^{1-p}$，需 $1-p>-1\Rightarrow p<2$；$x\to+\infty$：$\arctan x\sim\dfrac\pi2$，积分近 $x^{-(p+1)}$，需 $p+1>1\Rightarrow p>0$。故 $0<p<2$。
\par\noindent\textbf{12. $\dfrac12$}\par
$\lim\limits_{x\to0}\dfrac{\sin x-\ln(1+x)}{x\sin x}=\lim\limits_{x\to0}\dfrac{\left(x-\frac{x^3}{6}\right)-\left(x-\frac{x^2}{2}+\frac{x^3}{3}\right)}{x^2}=\lim\limits_{x\to0}\dfrac{\frac{x^2}{2}-\frac{x^3}{2}}{x^2}=\dfrac12$。
\par\noindent\textbf{13. $4$}\par
$F=x^2+2\sqrt3\,xy+y^2-1$，$F_x=2x+2\sqrt3\,y$，$F_y=2\sqrt3\,x+2y$。在 $(0,1)$：$F_x=2\sqrt3,\ F_y=2$，$y'=-\dfrac{F_x}{F_y}=-\sqrt3$。$y''=-\dfrac{F_{xx}F_y^2-2F_{xy}F_xF_y+F_{yy}F_x^2}{F_y^3}=\dfrac{16}{8}=2$。曲率 $\kappa=\dfrac{|y''|}{(1+y'^2)^{\frac32}}=\dfrac{2}{8}=\dfrac14$，半径 $R=4$。
\par\noindent\textbf{14. $-2\pi$}\par
$g'(x)=f_x(\ln x,\sin\pi x)\cdot\dfrac1x+f_y(\ln x,\sin\pi x)\cdot\pi\cos\pi x$。在 $x=1$：$(\ln x,\sin\pi x)=(0,0)$，$f_x(0,0)=\pi$，$f_y(0,0)=3$，故 $g'(1)=\pi+3\pi\cos\pi=\pi-3\pi=-2\pi$。
\par\small\color{gray}\emph{注：原答案页印作 $-2$，按上式应为 $-2\pi$（疑漏 $\pi$），据此修正。}\par
\par\noindent\textbf{15. $3\ln2-1$}\par
平均值 $\dfrac12\int_0^2\ln(2+x)\,\mathrm{d}x=\dfrac12\big[(x+2)\ln(x+2)-x\big]_0^2=\dfrac12\big[(4\ln4-4)-(2\ln2-2)\big]=\dfrac12(6\ln2-2)=3\ln2-1$。
\par\noindent\textbf{16. $2$}\par
$A=\begin{pmatrix}1&b&-1\\a+2&3&-3\end{pmatrix}$，$AA^{\mathrm{T}}=\begin{pmatrix}2+b^2&a+2+3b+3\\a+2+3b+3&(a+2)^2+18\end{pmatrix}$。规范形为 $y_1^2$（秩 $1$、一间号）要求 $AA^{\mathrm{T}}$ 秩 $1$ 且仅一个正特征值。由秩 $1$ 得 $a+2+3b+3=0$ 且 $2+b^2$ 与 $(a+2)^2+18$ 成比例，解得 $a+b=2$。
{\Large\bfseries 三、解答题（17～22）}\par\vspace{4pt}
\par\noindent\textbf{17. $4\sin\sqrt2-2\sqrt2$}\par
积分区域 $x^2+y^2\leqslant2$ 且 $y\geqslant|x|$。极坐标：$I=\displaystyle\int_{\frac\pi4}^{\frac{3\pi}4}\!\int_0^{\sqrt2}r^2\sin\theta\sin r\,\mathrm{d}r\mathrm{d}\theta=\left(\int_{\frac\pi4}^{\frac{3\pi}4}\sin\theta\,\mathrm{d}\theta\right)\left(\int_0^{\sqrt2}r^2\sin r\,\mathrm{d}r\right)=\sqrt2\left(2\sqrt2\sin\sqrt2-2\right)=4\sin\sqrt2-2\sqrt2$。
\par\noindent\textbf{18. $f'(x)=\begin{cases}\dfrac{3g(x^3)x^3-\int_0^{x^3}g(u)\,\mathrm{d}u}{x^2},&x\neq0\\ 0,&x=0\end{cases}$，且 $f'(x)$ 在 $x=0$ 处连续}\par
$x\neq0$ 时 $f(x)=\dfrac1x\int_0^{x^3}g(u)\,\mathrm{d}u$，故 $f'(x)=-\dfrac1{x^2}\int_0^{x^3}g(u)\,\mathrm{d}u+\dfrac1x\cdot3x^2g(x^3)=\dfrac{3x^3g(x^3)-\int_0^{x^3}g(u)\,\mathrm{d}u}{x^2}$。$f'(0)=\lim=\dfrac{3x^3g(x^3)-x^3g(\xi)}{x^2}\to0$（由 $g$ 连续），故 $f'(0)=0$ 且连续。
\par\noindent\textbf{19. 在 $(-2,0)$ 取极大值，极大值 $8\mathrm{e}^{-2}$}\par
$f_x=\mathrm{e}^x(2x^2+4x-y^2)$，$f_y=-2y\mathrm{e}^x$。由 $f_y=0$ 得 $y=0$，代入 $f_x=0$ 得 $2x(x+2)=0$，即 $(0,0),(-2,0)$。在 $(-2,0)$：$f_{xx}=-4\mathrm{e}^{-2}<0$，$f_{xy}=0$，$f_{yy}=-2\mathrm{e}^{-2}$，$f_{xx}f_{yy}-f_{xy}^2>0$，故为极大值，$f(-2,0)=8\mathrm{e}^{-2}$（$(0,0)$ 为鞍点）。
\par\noindent\textbf{20. $\dfrac{\pi^2}{6}-\dfrac{\sqrt3}{16}\pi$}\par
$y=\dfrac1{1+x^2}$ 拐点：$y''=\dfrac{-2+6x^2}{(1+x^2)^3}=0$，$x_0=\dfrac1{\sqrt3}$，$y_0=\dfrac34$。$V=\pi\left[\int_0^{1/\sqrt3}\left(\dfrac{3\sqrt3}{4}x\right)^2\mathrm{d}x+\int_{1/\sqrt3}^{+\infty}\dfrac{1}{(1+x^2)^2}\mathrm{d}x\right]=\pi\left[\dfrac{\sqrt3}{16}+\left(\dfrac\pi6-\dfrac{\sqrt3}{8}\right)\right]=\dfrac{\pi^2}{6}-\dfrac{\sqrt3}{16}\pi$。
\par\noindent\textbf{21. $y=-2x-\dfrac12x^2-4\ln|2-x|+11$}\par
令 $p=y'$，得 $x^2p'-2xp-p^2=0$。令 $u=\dfrac px$，化为 $xu'=u(1+u)$，积得 $\dfrac{p/x}{1+p/x}=Cx$。代入 $x=3,y'=-9$ 得 $C=\dfrac12$，故 $y'=\dfrac{x^2}{2-x}=-x-2-\dfrac4{x-2}$。积分 $y=-\dfrac{x^2}{2}-2x-4\ln|2-x|+C$，代入 $y(3)=\dfrac12$ 得 $C=11$。
\par\noindent\textbf{22.（1）$\alpha_1,\alpha_2$ 为极大线性无关组　（2）$H=\begin{pmatrix}1&0&-1&1\\0&1&1&-1\end{pmatrix}$，$A^{10}=\begin{pmatrix}1&-8&-9&9\\0&-1&-1&1\\-1&9&10&-10\\-1&7&8&-8\end{pmatrix}$}\par
（1）由 $\alpha_4=\alpha_1-\alpha_2$，$\alpha_3=\alpha_2-\alpha_1$，故 $\alpha_3,\alpha_4$ 均属 $\mathrm{span}\{\alpha_1,\alpha_2\}$；且 $\alpha_1,\alpha_2$ 线性无关，故其为极大线性无关组。
（2）由 $A=(\alpha_1,\alpha_2,\alpha_3,\alpha_4)$，$G=(\alpha_1,\alpha_2)$，$\alpha_3=-\alpha_1+\alpha_2$，$\alpha_4=\alpha_1-\alpha_2$，得 $H=\begin{pmatrix}1&0&-1&1\\0&1&1&-1\end{pmatrix}$。$HG=B=\begin{pmatrix}1&-1\\0&1\end{pmatrix}$，$B^{n-1}=\begin{pmatrix}1&-(n-1)\\0&1\end{pmatrix}$，$A^{10}=GB^9H$ 计算得如下矩阵
$A^{10}=\begin{pmatrix}1&-8&-9&9\\0&-1&-1&1\\-1&9&10&-10\\-1&7&8&-8\end{pmatrix}$。
\end{document}